%% ===================================================================
\section{Experimental Results}
\label{sec:eval}

\subsection{PolyBench/C (30 Kernels)}

\begin{table}[t]
  \caption{PolyBench/C results.  WA = with analysis, NA = no analysis.}
  \label{tab:results_1x}
  \centering\small
  \begin{tabular}{lcc}
    \toprule
    \textbf{Metric} & \textbf{WA} & \textbf{NA} \\
    \midrule
    Pass rate              & 30/30 (100\%) & 28/30 (93\%) \\
    Median speedup         & 2.07$\times$  & 1.06$\times$ \\
    Mean speedup           & 2.58$\times$  & 1.52$\times$ \\
    Kernels $>1\times$     & 20/28         & 14/28         \\
    \midrule
    \multicolumn{3}{l}{\emph{Head-to-head (28 common):}} \\
    WA wins ($>$10\% faster) & \multicolumn{2}{c}{12} \\
    NA wins ($>$10\% faster) & \multicolumn{2}{c}{6 (3 race conditions)}  \\
    Ties (within 10\%)       & \multicolumn{2}{c}{10}  \\
    \bottomrule
  \end{tabular}
\end{table}

Table~\ref{tab:results_1x} summarizes results at standard
(1$\times$) problem sizes.  All 30 kernels pass with analysis
guidance; without it, two stencil kernels fail from
parallelizing dimensions with carried dependences.
Compute-bound kernels show strong speedups
(\texttt{heat\_3d}: 10.27$\times$, \texttt{doitgen}: 7.40$\times$,
\texttt{covariance}: 5.40$\times$, \texttt{deriche}: 5.14$\times$).

In the head-to-head comparison, 3 of 6 apparent NA wins are \emph{false wins}: the unguided LLM parallelizes dimensions with carried dependences, producing code that runs faster but computes incorrect results. For example, in \texttt{lu}, the NA code launches one GPU thread per row (\texttt{grid=(N,)}), but row~$i$ depends on row~$k$ ($k<i$), causing output differences up to $10^{33}$.

\paragraph{Scaling.}
At 8$\times$ sizes ($N$: 60$\to$480, matrices: 120$\to$960), median speedup
rises to 30.73$\times$ over single-threaded C and 8.57$\times$ over 80-thread OpenMP,
confirming genuine advantage beyond CPU parallelism.

\subsection{Profiling-Guided Optimization}

\begin{table}[t]
  \caption{Representative profiling feedback improvements (PolyBench 1$\times$).
    ``Before'' = first correct kernel; ``After'' = best after NCU-guided iterations.}
  \label{tab:profiling}
  \centering\small
  \begin{tabular}{lccl}
    \toprule
    \textbf{Kernel} & \textbf{Before} & \textbf{After} & \textbf{Bottleneck} \\
    \midrule
    covariance  & 0.15$\times$ & 5.40$\times$ & latency \\
    deriche     & 0.11$\times$ & 5.14$\times$ & latency \\
    3mm         & 0.12$\times$ & 4.24$\times$ & latency \\
    correlation & 0.29$\times$ & 3.38$\times$ & latency \\
    adi         & 0.10$\times$ & 2.07$\times$ & latency \\
    gemver      & 0.16$\times$ & 1.94$\times$ & latency \\
    gemm        & 2.41$\times$ & 3.10$\times$ & latency \\
    fdtd\_2d    & 0.85$\times$ & 1.80$\times$ & latency \\
    durbin      & 0.30$\times$ & 1.47$\times$ & latency \\
    \bottomrule
  \end{tabular}
\end{table}

Profiling feedback improved 20 of 28 passing kernels (Table~\ref{tab:profiling} shows representative gains).
The largest improvements occur on latency-bound kernels where the initial implementation has insufficient parallelism: the NCU metrics reveal low SM and memory utilization ($<$20\%), and the LLM restructures the kernel to use more parallel resources.
For example, \texttt{covariance} improves from 0.15$\times$ to 5.40$\times$ (36$\times$ gain) as the LLM transforms a single-CTA sequential implementation into a multi-CTA parallel grid after receiving the latency-bound diagnosis.

Correctness is preserved across all profiling iterations: each optimized kernel is re-validated, and attempts that break correctness are automatically rejected.
The majority of optimization attempts do break correctness---the LLM frequently introduces race conditions when restructuring parallelism.
This confirms that single-shot prompting is insufficient for reliable optimization and motivates integration with search-based or RL-based methods that can explore many candidates.

\subsection{TSVC (151 Kernels)}

\sysname passes 150/151 TSVC kernels (99.3\%), 120 on first attempt.
At 32\,MB per array, LLM-generated Triton achieves 67$\times$ over single-threaded C,
10.7$\times$ over 80-thread OpenMP CPU, and \textbf{1.02$\times$ vs.\ OpenMP GPU offloading}
(\texttt{nvc~-mp=gpu}) on the same GPU---matching production compiler quality.

\subsection{Application Kernels}

\begin{table}[t]
  \caption{Application kernels: Triton vs.\ OMP~GPU offloading
    (\texttt{nvc}) on the same GPU.}
  \label{tab:realworld}
  \centering\small
  \begin{tabular}{llccc}
    \toprule
    \textbf{Source} & \textbf{Kernel} & \textbf{Tri/OMP\,GPU} & \textbf{WA/ST} & \textbf{NA/ST} \\
    \midrule
    Rodinia & SRAD       & 24.9$\times$  & 24.9$\times$ & 3.5$\times$ \\
    Rodinia & gaussian   & 249$\times$   & 3.4$\times$  & 3.6$\times$ \\
    miniMD  & LJ force   & 6.8$\times$   & 25.8$\times$ & 23.0$\times$ \\
    miniFE  & SpMV       & 3.1$\times$   & 9.0$\times$  & 7.3$\times$ \\
    Rodinia & lud        & ---$^a$       & 13.2$\times$ & 15.5$\times$ \\
    Rodinia & lavaMD     & ---$^a$       & 1.5$\times$  & 0.6$\times$ \\
    \bottomrule
    \multicolumn{5}{l}{\scriptsize $^a$OMP GPU offloading (\texttt{nvc}) fails to compile these kernels.}
  \end{tabular}
\end{table}

To evaluate generalization, we apply \sysname to 6 Rodinia and 2 ECP kernels (Table~\ref{tab:realworld}).
All 8 pass with analysis guidance.
On kernels with OMP~GPU baselines, Triton is 3--249$\times$ faster than \texttt{nvc}'s offloading, due to Triton's block-level model avoiding per-timestep launch overhead (gaussian: 511 timesteps $\times$ 3 launches = 1533 offloads).
Analysis guidance is most impactful on multi-phase kernels: SRAD improves from 3.5$\times$ (NA) to 24.9$\times$ (WA) via phase decomposition.
